% !TeX root = ../bd-screen.tex

\ifdefined\ollangid
\chapter{前言}

这是一本模态逻辑的入门教材。我在卡尔加里大学（University of Calgary）讲授哲学 579.2 课程（模态逻辑）时将其用作主要教材。它基于
\href{https://openlogicproject.org}{Open Logic Project} 的材料。

正文假定读者熟悉一些初等集合论和（命题）逻辑的基础。这些内容是我课程 Logic~II 的前置要求的组成部分。该课程的教材
\href{https://slc.openlogicproject.org}{\emph{Sets, Logic, Computation}} 同样基于 OLP，因此可免费获取。不过，所需的材料也作为附录收录在本书中。我每次讲授这些内容时，都会布置这些附录供背景阅读。

\Cref{nml:::part} 部分最初部分基于 Aldo Antonelli 关于「基本模态逻辑的经典对应理论」的讲稿，这些讲稿是他于 2015 年不幸去世前贡献给 OLP 的。我对原稿做了大幅修订与扩充，例如，关于框架可定义性和表列的材料即为新增内容。
\else
\chapter{Preface}

This is an introductory textbook on modal logic. I use it as the main
text when I teach Philosophy~579.2 (Modal Logic) at the University of
Calgary. It is based on material from the
\href{https://openlogicproject.org}{Open Logic Project}.

The main text assumes familiarity with some elementary set theory and
the basics of (propositional) logic. This material is part of a
prerequisite for my course, Logic~II. The textbook for that course,
\href{https://slc.openlogicproject.org}{\emph{Sets, Logic,
Computation}}, is also based on the OLP, and so is available for free.
The required material is included as appendices in this book,
however. I assign these appendices for background reading whenever I
teach the material.

\Cref{nml:::part} is originally based in part on Aldo Antonelli's
lecture notes on ``Classical Correspondence Theory for Basic Modal
Logic,'' which he contributed to the OLP before his untimely death in
2015. I heavily revised and expanded these notes, e.g., the material
on frame definability and tableaux is new.
\fi
